%! TEX root = thesis.tex

\chapter{Conclusion}
\label{sec:conclusion}

We have presented a simulation framework and adaptive grid-search optimizer
for carbon-aware LLM training via hysteresis-based pausing.  The framework
replays historical grid carbon intensity data through a discrete-event
simulator, and the optimizer efficiently searches the three-dimensional policy
space of pause threshold, resume threshold, and start date.  Across two
state-of-the-art models (DeepSeek~V3 and Kimi~K2) and five grid regions with
diverse variability profiles, we establish systematic Pareto frontiers for
this problem.

The central finding is that hysteresis provides diminishing returns: the
optimal policy uses a narrow or zero margin ($\Delta_\theta \leq
16$~gCO$_2$eq/kWh) in every region and model.  Wide hysteresis margins waste
time and checkpoint overhead without proportional savings.  Carbon-aware
pausing can reduce operational emissions by up to 43\% in high-variability
grids at the cost of up to 200\% wall-clock overhead.  In low-variance grids,
savings fall below 10\%, and alternative mitigation strategies such as
geo-distributed training should be considered.

Our implementation as a fully client-side browser application demonstrates
that carbon-aware training simulation and optimization can be made accessible
to researchers and practitioners without specialized infrastructure.  The
application provides an interactive live simulator, a batch run dashboard,
and an adaptive optimization page with real-time 3D visualization of the
policy space.

\section{Limitations}

Our simulation assumes perfect foresight: policies are evaluated on historical
data rather than forecasted future intensity.  A real deployment would require
forecasting, introducing prediction error.  We use a fixed checkpoint overhead
derived from a 7B model; real overhead depends on model size and storage
bandwidth and may scale with the number of model parameters.  GPU idle power
is an estimate, and actual values vary across hardware generations and data
center configurations.  Finally, this work considers single-site training;
geo-distributed training could compound savings by adding spatial flexibility.

\section{Future Work}

Three directions are particularly promising for future work:
\begin{enumerate}
  \item \textbf{Forecast-based policies.}  Replace historical replay with
    predictive carbon intensity forecasts (e.g.\ using ECMWF or local grid
    operator data) to enable real-time deployment.
  \item \textbf{Multi-job scheduling.}  Extend the framework to optimize across
    multiple training jobs sharing a GPU cluster, balancing carbon savings with
    fairness and throughput.
  \item \textbf{Embodied carbon awareness.}  Incorporate manufacturing
    emissions to ensure that operational savings from pausing are not
    outweighed by amortized hardware production costs.
\end{enumerate}

All simulation code and data are available as open source at
\url{https://github.com/404simon/TheGreenEpoch}.
